\documentclass[a4paper,12pt]{article}
\usepackage[T2A]{fontenc}
\usepackage[utf8]{inputenc}
\usepackage[russian]{babel}
\usepackage{amsmath,amssymb,mathtools}
\usepackage{icomma}
\usepackage[left=22mm,right=22mm,top=23mm,bottom=23mm,headheight=26pt]{geometry}
\usepackage{microtype}
\usepackage{xcolor}
\usepackage{tikz}
\usetikzlibrary{calc,arrows.meta,positioning,shapes.geometric}
\usepackage{tkz-euclide}
\usepackage{booktabs,tabularx}
\usepackage{enumitem}
\usepackage{fancyhdr}
\usepackage{hyperref}
\usepackage{parskip}
\usepackage{setspace}
\definecolor{accent}{HTML}{007F84}
\definecolor{darkaccent}{HTML}{075459}
\definecolor{textgray}{HTML}{34383A}
\definecolor{softgray}{HTML}{697174}
\definecolor{softline}{HTML}{B6C9C9}
\hypersetup{colorlinks=true,linkcolor=darkaccent,pdfauthor={Лёвин Артём Александрович},pdftitle={Чек-лист: правильные пирамиды},pdfsubject={Профильная математика. Николь Саркисьянц}}
\setlength{\parindent}{0pt}
\setlength{\parskip}{6pt plus 1pt}
\setstretch{1.16}
\setlist[itemize]{leftmargin=6mm,itemsep=3pt,topsep=3pt}
\setlist[enumerate]{leftmargin=7mm,itemsep=8pt,topsep=4pt}
\pagestyle{fancy}
\fancyhf{}
\fancyhead[L]{\small\color{textgray}Правильные пирамиды}
\fancyhead[R]{\small\color{textgray}Николь Саркисьянц}
\fancyfoot[C]{\small\color{softgray}\thepage}
\renewcommand{\headrulewidth}{0.35pt}
\renewcommand{\headrule}{\hbox to\headwidth{\color{accent}\leaders\hrule height \headrulewidth\hfill}}
\newcommand{\sect}[1]{\section{#1}\vspace{-1.5mm}\noindent\color{accent}\rule{\linewidth}{0.45pt}\color{textgray}\vspace{2mm}}
\newcommand{\keyidea}[1]{\par\textcolor{darkaccent}{\textbf{Ключевая идея.}} #1\par}
\newcommand{\warning}[1]{\par\textcolor{darkaccent}{\textbf{Проверяйте.}} #1\par}
\newcommand{\ex}[1]{\par\medskip\noindent\textcolor{darkaccent}{\textbf{#1}}\par}
\newcommand{\answer}[1]{\textcolor{darkaccent}{\(\displaystyle #1\)}}
\newcommand{\cb}{\(\square\)}
\begin{document}
\pagenumbering{gobble}
\color{textgray}
\thispagestyle{empty}
\begin{titlepage}
\begin{center}
\vspace*{22mm}
{\large\color{darkaccent} МАТЕМАТИКА}\par
\vspace{23mm}
{\LARGE Чек-лист}\par
\vspace{13mm}
{\Large\color{accent} Правильные пирамиды:\par}
{\Large\color{accent} объём, площадь поверхности\par}
{\Large\color{accent} и подобие}\par
\vspace{21mm}
{\normalsize Подготовка к ЕГЭ по профильной математике}\par
\vfill
{\large Николь Саркисьянц}\par
\vspace{12mm}
{\normalsize 19 ноября 2026 г.}\par
\vspace{18mm}
{\small Лёвин Артём Александрович\par эксклюзивно для Николь Саркисьянц}\par
\vspace{21mm}
\begin{tikzpicture}
\draw[accent,line width=.75pt] (-4,0) .. controls (-3,0.62) and (-2,0.62) .. (-1,0)
.. controls (0,-.36) and (.7,-.25) .. (1.2,0)
.. controls (2.1,.38) and (3.1,.5) .. (4,0);
\end{tikzpicture}
\end{center}
\end{titlepage}
\pagenumbering{arabic}
\renewcommand{\thesection}{\arabic{section}.}

\sect{С чего начинается задача о пирамиде}
\textbf{Пирамида} имеет одно основание, а остальные грани --- треугольники. У \textbf{правильной пирамиды} основание --- правильный многоугольник, а вершина расположена над его центром. Боковые рёбра равны между собой.

Обозначения, которые будем использовать:
\begin{itemize}
\item \(S\) --- вершина пирамиды; \(O\) --- центр основания;
\item \(H=SO\) --- \textbf{высота пирамиды}, перпендикулярная плоскости основания;
\item \(M\) --- середина стороны основания; \(h=SM\) --- \textbf{апофема}, высота боковой грани;
\item \(a\) --- сторона основания, \(P\) --- периметр основания, \(B\) --- его площадь.
\end{itemize}

\begin{center}
\begin{minipage}[c]{.54\linewidth}\centering
\begin{tikzpicture}[scale=.73]
\tkzDefPoint(-2,-.8){A}
\tkzDefPoint(1,-1.2){B}
\tkzDefPoint(3,.2){C}
\tkzDefPoint(0,.6){D}
\tkzDefPoint(.5,-.3){O}
\tkzDefPoint(.5,2.65){S}
\tkzDefPoint(2,-.5){M}
\tkzDrawSegments[gray!70,dashed,line width=.6pt](A,D D,C S,D A,C B,D S,O O,M)
\tkzDrawSegments[darkaccent,line width=1pt](A,B B,C S,A S,B S,C)
\tkzDrawSegments[accent,line width=1pt](S,M)
\tkzLabelPoints[below left](A)
\tkzLabelPoints[below](B)
\tkzLabelPoints[right](C)
\tkzLabelPoints[above left](D)
\tkzLabelPoints[above](S)
\tkzLabelPoints[below left](O)
\tkzLabelPoints[below right](M)
\end{tikzpicture}
\end{minipage}\hfill
\begin{minipage}[c]{.43\linewidth}\centering
\begin{tikzpicture}[scale=.85]
\tkzDefPoint(0,0){O}
\tkzDefPoint(3.2,0){M}
\tkzDefPoint(0,3.6){S}
\tkzDrawSegments[darkaccent,line width=1pt](S,O O,M S,M)
\tkzMarkRightAngle[accent,size=.28](S,O,M)
\tkzLabelPoints[below left](O)
\tkzLabelPoints[below right](M)
\tkzLabelPoints[above](S)
\tkzLabelSegment[left](S,O){$H$}
\tkzLabelSegment[below](O,M){$a/2$}
\tkzLabelSegment[right](S,M){$h$}
\end{tikzpicture}
\end{minipage}
\end{center}
\smallskip
\textit{Слева --- схема правильной четырёхугольной пирамиды, справа --- настоящий прямоугольный треугольник \(SOM\). Внутренние отрезки на пространственном чертеже показаны пунктиром.}

\keyidea{Высота \(SO\) и апофема \(SM\) --- разные отрезки. Площадь боковой поверхности вычисляют через апофему; объём --- через высоту пирамиды.}

\sect{Формулы объёма и площади}
Для \textbf{любой} пирамиды:
\[
\boxed{V=\frac13 BH},\qquad
\boxed{S_{\text{полн}}=S_{\text{бок}}+B}.
\]
Для \textbf{правильной} пирамиды с апофемой \(h\):
\[
\boxed{S_{\text{бок}}=\frac12 Ph}.
\]
Здесь \(P\) --- периметр основания. Объяснение: все боковые грани имеют одинаковую высоту \(h\), поэтому сумма их площадей равна половине произведения периметра на \(h\).

Если основание --- квадрат со стороной \(a\), то
\[
B=a^2=\frac{d^2}{2},\quad P=4a,\quad
\boxed{S_{\text{бок}}=2ah},\quad
\boxed{S_{\text{полн}}=a^2+2ah},
\]
где \(d=a\sqrt2\) --- диагональ квадрата.

\ex{Пример 1. Объём через диагональ и боковое ребро}
Диагональ квадратного основания равна \(8\), боковое ребро --- \(5\). Найдём объём.
\[
B=\frac{8^2}{2}=32,\qquad
OA=\frac82=4,
\]
\[
H=\sqrt{SA^2-OA^2}=\sqrt{5^2-4^2}=3,
\]
\[
\boxed{V=\frac13\cdot32\cdot3=32}.
\]
\warning{Здесь \(OA\) --- половина \emph{диагонали}. Отрезок \(OM\) --- половина \emph{стороны}. Их нельзя путать.}

\sect{Как находить апофему}
Есть два полезных прямоугольных треугольника:
\[
\underbrace{SA^2=H^2+OA^2}_{\triangle SOA},
\qquad
\underbrace{h^2=H^2+OM^2}_{\triangle SOM}.
\]
Для квадратного основания \(OM=a/2\). Для боковой грани \(SAB\), где \(M\) --- середина \(AB\), верно \(AM=a/2\), поэтому
\[
\boxed{h=\sqrt{SA^2-\left(\frac a2\right)^2}}
\quad\text{или}\quad
\boxed{h=\sqrt{H^2+\left(\frac a2\right)^2}}.
\]

\ex{Пример 2. Боковая и полная поверхность}
Сторона квадратного основания равна \(6\), каждое боковое ребро --- \(5\). В боковой грани имеем
\[
AM=3,\qquad h=\sqrt{5^2-3^2}=4.
\]
Одна боковая грань: \(S_{\triangle SAB}=\tfrac12\cdot6\cdot4=12\).
Всего боковых граней \textbf{четыре}, следовательно,
\[
S_{\text{бок}}=4\cdot12=48,\qquad
\boxed{S_{\text{полн}}=48+6^2=84}.
\]

\ex{Пример 3. Высота известна, апофему ищем}
Если \(a=6\) и \(H=4\), то \(OM=3\), откуда
\[
h=\sqrt{4^2+3^2}=5,\qquad
S_{\text{бок}}=\frac{4\cdot6\cdot5}{2}=60,
\quad S_{\text{полн}}=96.
\]
\keyidea{Алгоритм для площади поверхности: определите, какое именно \(S\) требуется, найдите \(h\), вычислите \(S_{\text{бок}}\), затем при необходимости прибавьте \(B\).}

\sect{Правильный шестиугольник: что важно знать}
Правильный шестиугольник делится отрезками от центра к вершинам на \textbf{шесть равносторонних треугольников}. Каждый центральный угол равен \(360^\circ/6=60^\circ\).

\begin{center}
\begin{tikzpicture}[scale=.64]
\tkzDefPoint(0,0){O}
\tkzDefPoint(3,0){A}
\tkzDefPoint(1.5,2.598076){B}
\tkzDefPoint(-1.5,2.598076){C}
\tkzDefPoint(-3,0){D}
\tkzDefPoint(-1.5,-2.598076){E}
\tkzDefPoint(1.5,-2.598076){F}
\tkzDrawPolygon[darkaccent,line width=.95pt](A,B,C,D,E,F)
\tkzDrawSegments[gray!65,dashed](O,A O,B O,C O,D O,E O,F)
\tkzDrawSegments[accent,line width=1.15pt](A,C)
\tkzMarkAngle[accent,size=.85](A,O,B)
\tkzLabelAngle[font=\scriptsize,pos=1.2](A,O,B){$60^\circ$}
\tkzLabelPoints[right](A)
\tkzLabelPoints[above right](B)
\tkzLabelPoints[above left](C)
\tkzLabelPoints[left](D)
\tkzLabelPoints[below left](E)
\tkzLabelPoints[below right](F)
\tkzLabelPoints[below](O)
\tkzLabelSegment[right](A,B){$a$}
\tkzLabelSegment[above right](A,C){$a\sqrt3$}
\end{tikzpicture}
\end{center}
Площадь одного равностороннего треугольника \(\sqrt3 a^2/4\), поэтому
\[
\boxed{B_6=6\cdot\frac{\sqrt3a^2}{4}=\frac{3\sqrt3}{2}a^2}.
\]
Ещё три свойства, которые помогают в задачах:
\[
\boxed{OA=a}\quad\text{(радиус описанной окружности)},
\qquad \boxed{AC=a\sqrt3},\qquad \boxed{AD=2a}.
\]
Здесь \(AC\) --- диагональ через одну вершину, \(AD\) --- диагональ через центр.
Каждый внутренний угол правильного шестиугольника равен \(120^\circ\).

\ex{Пример 4. Боковое ребро шестиугольной пирамиды}
Дано: \(a=1\), \(V=6\). Площадь основания и высота:
\[
B_6=\frac{3\sqrt3}{2},\qquad
H=\frac{3V}{B_6}=\frac{18}{3\sqrt3/2}=\frac{12}{\sqrt3}=4\sqrt3.
\]
Так как \(OA=a=1\), получаем
\[
SA=\sqrt{H^2+OA^2}=\sqrt{(4\sqrt3)^2+1^2}=\boxed{7}.
\]

\sect{Угол между боковой гранью и основанием}
Пусть \(M\) --- середина стороны квадратного основания, \(O\) --- центр основания. Угол между боковой гранью и основанием измеряют \textbf{углом между} апофемой \(SM\) \textbf{и} её проекцией \(OM\) на основание:
\[
\angle SMO=\varphi,\qquad \tg\varphi=\frac{SO}{OM}=\frac{H}{a/2}.
\]
Если \(\varphi=45^\circ\), то \(H=OM=a/2\).

\ex{Пример 5. Объём при угле \(45^\circ\)}
Если \(a=8\), \(\varphi=45^\circ\), то \(H=4\), \(B=64\) и
\[
\boxed{V=\frac13\cdot64\cdot4=\frac{256}{3}}.
\]
\warning{На чертеже угол между плоскостями показывают в сечении, перпендикулярном их линии пересечения. Угол при вершине \(S\) в боковом треугольнике для этого не подходит.}

\sect{Правильный тетраэдр и масштабирование}
\textbf{Правильный тетраэдр} состоит из четырёх равносторонних треугольников. При ребре \(a\):
\[
\boxed{S_{\text{полн}}=4\cdot\frac{\sqrt3a^2}{4}=\sqrt3a^2},
\quad
H=a\sqrt{\frac23},\quad
\boxed{V=\frac{\sqrt2}{12}a^3}.
\]
Формула высоты следует из треугольника, соединяющего вершину тетраэдра, центр равностороннего основания и вершину основания: \(AO=a/\sqrt3\) и
\[
H^2=a^2-\frac{a^2}{3}=\frac{2a^2}{3}.
\]

При \textbf{одновременном} умножении всех линейных размеров тела на \(k>0\):
\[
\boxed{\frac{S_2}{S_1}=k^2},\qquad
\boxed{\frac{V_2}{V_1}=k^3}.
\]
Например, если каждое ребро тетраэдра увеличили вдвое, площадь поверхности возрастёт в \(4\) раза, а объём --- в \(8\) раз.

\warning{Если сторону основания умножили на \(2\), а высоту --- на \(3\), это \emph{не} равномерное масштабирование. По формуле \(V=BH/3\) объём пирамиды увеличится в \(2^2\cdot3=12\) раз.}

\sect{Проверка перед записью ответа}
\begin{tabularx}{\linewidth}{@{}p{9mm}X@{}}
\cb& Я обозначил вершину, центр основания, высоту и апофему.\\[2pt]
\cb& Я отличаю боковое ребро от высоты боковой грани.\\[2pt]
\cb& В задаче на площадь определил: одна грань, вся боковая или полная поверхность.\\[2pt]
\cb& Для неизвестного отрезка выбрал подходящий прямоугольный треугольник.\\[2pt]
\cb& При умножении длин на \(k\) проверил, все ли размеры меняются одинаково.\\[2pt]
\cb& Указал точный ответ и убедился, что площадь/объём положительны.\\
\end{tabularx}

\clearpage
\thispagestyle{fancy}
\begin{center}{\Large\color{darkaccent} Алгоритм решения задачи о правильной пирамиде}\end{center}
\vspace{10mm}
\begin{center}
\begin{tikzpicture}[
>=Stealth, line width=.75pt,
term/.style={ellipse,draw=darkaccent,minimum width=6.0cm,minimum height=1cm,align=center,font=\small},
action/.style={rounded corners=2pt,draw=darkaccent,minimum width=6.6cm,text width=6.15cm,minimum height=1.2cm,align=center,font=\small,inner sep=7pt},
decision/.style={diamond,aspect=2.8,draw=accent,minimum width=5.3cm,minimum height=1.85cm,text width=3.65cm,align=center,font=\small,inner sep=1pt},
side/.style={rounded corners=2pt,draw=darkaccent,minimum width=4.5cm,text width=4.2cm,minimum height=1.65cm,align=center,font=\small,inner sep=6pt}]
\node[term] (start) at (-2,0) {Прочитайте условие и сделайте чертёж};
\node[action,below=1.05cm of start] (data) {Отметьте \(O\), высоту \(H\), апофему \(h\), известные рёбра и стороны};
\node[action,below=.95cm of data] (formula) {Выберите формулу: \(V=BH/3\), \(S_{\text{бок}}=Ph/2\) или \(S_{\text{полн}}=S_{\text{бок}}+B\)};
\node[decision,below=1.1cm of formula] (check) {Все нужные величины известны?};
\node[side,right=1.15cm of check] (find) {Найдите недостающее через площадь основания и прямоугольный треугольник};
\node[action,below=1.2cm of check] (calc) {Подставьте числа, выполните вычисления};
\node[action,below=.95cm of calc] (verify) {Проверьте формулу, тип площади и единицы измерения};
\node[term,below=.95cm of verify] (finish) {Запишите ответ};
\draw[->] (start)--(data);
\draw[->] (data)--(formula);
\draw[->] (formula)--(check);
\draw[->] (check)--node[left,font=\small]{да}(calc);
\draw[->] (check)--node[above,font=\small]{нет}(find);
\draw[->] (find.north) |- ([yshift=.47cm]formula.east) -- (formula.east);
\draw[->] (calc)--(verify);
\draw[->] (verify)--(finish);
\end{tikzpicture}
\end{center}
\clearpage

\begin{center}{\Large\color{darkaccent} Тренировочный блок}\par
\vspace{2mm}{\normalsize Путь А --- 10 задач от простого к сложному}\end{center}
\smallskip
Решайте последовательно. В задачах с площадью поверхности обязательно указывайте, какая именно площадь найдена.
\begin{enumerate}
\item Сторона основания правильной четырёхугольной пирамиды равна \(6\). Найдите площадь основания.
\item Диагональ основания правильной четырёхугольной пирамиды равна \(12\). Найдите площадь основания.
\item Сторона квадратного основания равна \(8\), боковое ребро --- \(5\). Найдите апофему пирамиды.
\item Сторона квадратного основания равна \(10\), боковое ребро --- \(13\). Найдите \emph{боковую} и \emph{полную} площади поверхности.
\item У правильной четырёхугольной пирамиды \(a=8\), высота \(H=3\). Найдите апофему и полную площадь поверхности.
\item Диагональ квадратного основания равна \(10\), боковое ребро --- \(\sqrt{34}\). Найдите объём пирамиды.
\item Сторона правильного шестиугольника равна \(4\). Найдите его площадь и длину диагонали, соединяющей вершины через одну.
\item У правильной шестиугольной пирамиды сторона основания \(2\), боковое ребро \(4\). Найдите объём.
\item Сторона основания правильной четырёхугольной пирамиды равна \(12\), угол между боковой гранью и плоскостью основания равен \(45^\circ\). Найдите объём и полную площадь поверхности.
\item В правильном тетраэдре все рёбра увеличили в \(2\) раза. Во сколько раз изменились площадь полной поверхности и объём? Затем рассмотрите \emph{другую} пирамиду: сторону её квадратного основания увеличили в \(2\) раза, а высоту --- в \(3\) раза. Во сколько раз изменился её объём?
\end{enumerate}

\clearpage
\begin{center}{\Large\color{darkaccent} Ответы}\end{center}
\vspace{6mm}
\renewcommand{\arraystretch}{1.65}
\begin{tabularx}{\linewidth}{@{}>{\centering\arraybackslash}p{11mm}X@{}}
\toprule
№ & Ответ \\
\midrule
1 & \(36\) \\
2 & \(72\) \\
3 & \(3\) \\
4 & \(S_{\text{бок}}=240,\quad S_{\text{полн}}=340\) \\
5 & \(h=5,\quad S_{\text{полн}}=144\) \\
6 & \(50\) \\
7 & \(24\sqrt3,\quad 4\sqrt3\) \\
8 & \(12\) \\
9 & \(V=288,\quad S_{\text{полн}}=144+144\sqrt2\) \\
10 & Тетраэдр: площадь в \(4\) раза, объём в \(8\) раз; другая пирамида: объём в \(12\) раз. \\
\bottomrule
\end{tabularx}
\vspace{16mm}
\begin{center}\textcolor{darkaccent}{\large Итоговая памятка}\end{center}
\renewcommand{\arraystretch}{1.33}
\begin{tabularx}{\linewidth}{@{}p{47mm}>{\raggedright\arraybackslash}X@{}}
\toprule
Сигнал в условии & Первое полезное действие \\
\midrule
«Объём пирамиды» & \(V=BH/3\); найдите \(B\) и \(H\).\\
«Полная поверхность» & Найдите \(Ph/2\) и прибавьте площадь основания.\\
«Известно боковое ребро» & Постройте треугольник с этим ребром в качестве гипотенузы.\\
«Шестиугольник» & \(B=3\sqrt3a^2/2\), радиус описанной окружности равен \(a\).\\
«Угол между гранью и основанием» & Постройте сечение через высоту пирамиды и апофему.\\
«Все размеры увеличили в \(k\) раз» & Площадь умножается на \(k^2\), объём --- на \(k^3\).\\
\bottomrule
\end{tabularx}
\vfill
\begin{center}\small\color{softgray}Лёвин Артём Александрович эксклюзивно для Николь Саркисьянц\end{center}
\end{document}
